\documentclass[10pt,a4paper]{amsart}
\usepackage{amsmath,amssymb,amsthm,graphicx}
\usepackage{verbatim}
\newenvironment{reference}{\comment}{\endcomment}
%\newenvironment{reference}{}{}
\newenvironment{slogan}{\comment}{\endcomment}
\newenvironment{history}{\comment}{\endcomment}

% For commutative diagrams we use Xy-pic
\usepackage[all]{xy}

% We use 2cell for 2-commutative diagrams.
\xyoption{2cell}
% Theorem environments.
%
\theoremstyle{plain}
\newtheorem{theorem}[subsection]{Theorem}
\newtheorem{proposition}[subsection]{Proposition}
\newtheorem{lemma}[subsection]{Lemma}

\theoremstyle{definition}
\newtheorem{definition}[subsection]{Definition}
\newtheorem{example}[subsection]{Example}
\newtheorem{exercise}[subsection]{Exercise}
\newtheorem{situation}[subsection]{Situation}

\theoremstyle{remark}
\newtheorem{remark}[subsection]{Remark}
\newtheorem{remarks}[subsection]{Remarks}

\numberwithin{equation}{subsection}

% Macros
%
\def\lim{\mathop{\mathrm{lim}}\nolimits}
\def\colim{\mathop{\mathrm{colim}}\nolimits}
\def\Spec{\mathop{\mathrm{Spec}}}
\def\Hom{\mathop{\mathrm{Hom}}\nolimits}
\def\Ext{\mathop{\mathrm{Ext}}\nolimits}
\def\SheafHom{\mathop{\mathcal{H}\!\mathit{om}}\nolimits}
\def\SheafExt{\mathop{\mathcal{E}\!\mathit{xt}}\nolimits}
\def\Sch{\mathit{Sch}}
\def\Mor{\mathop{\mathrm{Mor}}\nolimits}
\def\Ob{\mathop{\mathrm{Ob}}\nolimits}
\def\Sh{\mathop{\mathit{Sh}}\nolimits}
\def\NL{\mathop{N\!L}\nolimits}
\def\CH{\mathop{\mathrm{CH}}\nolimits}
\def\proetale{{pro\text{-}\acute{e}tale}}
\def\etale{{\acute{e}tale}}
\def\QCoh{\mathit{QCoh}}
\def\Ker{\mathop{\mathrm{Ker}}}
\def\Im{\mathop{\mathrm{Im}}}
\def\Coker{\mathop{\mathrm{Coker}}}
\def\Coim{\mathop{\mathrm{Coim}}}

% Boxtimes
%
\DeclareMathSymbol{\boxtimes}{\mathbin}{AMSa}{"02}

%
% Macros for moduli stacks/spaces
%
\def\QCohstack{\mathcal{QC}\!\mathit{oh}}
\def\Cohstack{\mathcal{C}\!\mathit{oh}}
\def\Spacesstack{\mathcal{S}\!\mathit{paces}}
\def\Quotfunctor{\mathrm{Quot}}
\def\Hilbfunctor{\mathrm{Hilb}}
\def\Curvesstack{\mathcal{C}\!\mathit{urves}}
\def\Polarizedstack{\mathcal{P}\!\mathit{olarized}}
\def\Complexesstack{\mathcal{C}\!\mathit{omplexes}}
% \Pic is the operator that assigns to X its picard group, usage \Pic(X)
% \Picardstack_{X/B} denotes the Picard stack of X over B
% \Picardfunctor_{X/B} denotes the Picard functor of X over B
\def\Pic{\mathop{\mathrm{Pic}}\nolimits}
\def\Picardstack{\mathcal{P}\!\mathit{ic}}
\def\Picardfunctor{\mathrm{Pic}}
\def\Deformationcategory{\mathcal{D}\!\mathit{ef}}
\usepackage{cleveref}
\crefname{originallemma}{original lemma}{original lemmas}
\Crefname{originallemma}{Original lemma}{Original lemmas}
\makeatletter
\newlabel{lemma-equivalent-definitions}{{0A86}{source}{}{sourcepointer-1}{}}
\newlabel{lemma-equivalent-definitions@cref}{{[originallemma][0][]0A86}{[1][0][]source}{}{}{}}
\makeatother
\crefname{originallemma}{original lemma}{original lemmas}
\Crefname{originallemma}{Original lemma}{Original lemmas}
\makeatletter
\newlabel{lemma-affine-duality}{{0A88}{source}{}{sourcepointer-2}{}}
\newlabel{lemma-affine-duality@cref}{{[originallemma][0][]0A88}{[1][0][]source}{}{}{}}
\makeatother
\crefname{originallemma}{original lemma}{original lemmas}
\Crefname{originallemma}{Original lemma}{Original lemmas}
\makeatletter
\newlabel{cohomology-lemma-internal-hom-evaluate}{{0A8U}{source}{}{sourcepointer-3}{}}
\newlabel{cohomology-lemma-internal-hom-evaluate@cref}{{[originallemma][0][]0A8U}{[1][0][]source}{}{}{}}
\makeatother
\crefname{originallemma}{original lemma}{original lemmas}
\Crefname{originallemma}{Original lemma}{Original lemmas}
\makeatletter
\newlabel{dualizing-lemma-dualizing}{{0A7C}{source}{}{sourcepointer-4}{}}
\newlabel{dualizing-lemma-dualizing@cref}{{[originallemma][0][]0A7C}{[1][0][]source}{}{}{}}
\makeatother
\crefname{originallemma}{original lemma}{original lemmas}
\Crefname{originallemma}{Original lemma}{Original lemmas}
\makeatletter
\newlabel{lemma-internal-hom-evaluate-isom}{{0G4I}{source}{}{sourcepointer-5}{}}
\newlabel{lemma-internal-hom-evaluate-isom@cref}{{[originallemma][0][]0G4I}{[1][0][]source}{}{}{}}
\makeatother
\title{A dualizing complex on a locally Noetherian scheme induces an anti-equivalence of the bounded coherent derived category}
\author{The Stacks Project Authors}
\date{stacks-project:duality.tex:local-snapshot-20261012; GNU FDL 1.2 or later}
\begin{document}
\maketitle
\noindent\textit{Editorial note.} All delivered mathematical text is the author's own, copied from the pinned Stacks Project chapter file duality.tex; no statement or proof step was rewritten, reordered or removed.\par\medskip
\noindent\textit{Editorial note.} Preamble provenance: the pinned chapter file opens with an input of the shared preamble and carries no preamble text of its own, so the project's single Stacks-wide preamble staged at sources/stacks-duality/preamble.tex is excerpted verbatim under the recipe fixed by LANE\_BRIEF 22.1, as three spans, lines 8-13 (verbatim package and the reference/slogan/history environments), lines 14-18 (Xy-pic setup) and lines 154-217 (the native theorem declarations and mathematical macros). Nothing was added, removed or redefined; the wrapper class and the amsmath/amssymb/amsthm/graphicx base packages are supplied by the assembler.\par\medskip
\noindent\textit{Editorial note.} The counted statement and proof contain no cross reference, no citation and no included file, so the delivered body is self contained.\par\medskip
\noindent\textit{Editorial note.} Canonical location of the counted proposition: Stacks tag 0A89 (sources/stacks-duality/pinned-original-source.tex line 472).\par\medskip
\noindent\textit{Allowed standard prerequisites (editorial).}\begin{itemize}
\item Comparison of the scheme level and affine level dualizing functors. Same chapter, Lemmas labelled lemma-equivalent-definitions and lemma-affine-duality, cited at proof lines 491-506. Shows the functor D is induced by the algebraic dualizing functor on affines.
\item The algebraic dualizing functor is an anti-equivalence with id -> D o D. Dualizing Complexes, Lemma labelled dualizing-lemma-dualizing, cited at proof lines 518-522 and 524-526. Supplies the isomorphism and the boundedness statements.
\item The evaluation map between internal Homs. Cohomology, Lemma labelled cohomology-lemma-internal-hom-evaluate, cited at proof lines 512-516. Provides the second arrow of the displayed composition.
\item The evaluation isomorphism alternative. Same chapter, Lemma labelled lemma-internal-hom-evaluate-isom part (2), cited in the footnote at proof lines 520-524. Gives the alternative route to the isomorphism.
\end{itemize}\medskip
\ifdefined\hypertarget\hypertarget{sourcepointer-1}{}\fi
\noindent\textit{Source reference.} Original lemma 0A86: sources/stacks-duality/pinned-original-source.tex, source lines 308--308. Same-chapter pointer into the author environment lemma at source line 308; invoked at proof line 493 of the counted proof. It is retained only as a navigation pointer; every load-bearing step is inside the delivered span.\par\medskip
\ifdefined\hypertarget\hypertarget{sourcepointer-2}{}\fi
\noindent\textit{Source reference.} Original lemma 0A88: sources/stacks-duality/pinned-original-source.tex, source lines 366--366. Same-chapter pointer into the author environment lemma at source line 366; invoked at proof line 494 of the counted proof. It is retained only as a navigation pointer; every load-bearing step is inside the delivered span.\par\medskip
\ifdefined\hypertarget\hypertarget{sourcepointer-3}{}\fi
\noindent\textit{Source reference.} Original lemma 0A8U: sources/stacks-chow/tags.txt, source lines 12927--12927. Cross-chapter standard prerequisite located in the project tag list staged in the local snapshot; invoked at proof line 512 of the counted proof. It is retained only as a navigation pointer; every load-bearing step is inside the delivered span.\par\medskip
\ifdefined\hypertarget\hypertarget{sourcepointer-4}{}\fi
\noindent\textit{Source reference.} Original lemma 0A7C: sources/stacks-chow/tags.txt, source lines 12875--12875. Cross-chapter standard prerequisite located in the project tag list staged in the local snapshot; invoked at proof line 515 of the counted proof. It is retained only as a navigation pointer; every load-bearing step is inside the delivered span.\par\medskip
\ifdefined\hypertarget\hypertarget{sourcepointer-5}{}\fi
\noindent\textit{Source reference.} Original lemma 0G4I: sources/stacks-duality/pinned-original-source.tex, source lines 415--415. Same-chapter pointer into the author environment lemma at source line 415; invoked at proof line 519 of the counted proof. It is retained only as a navigation pointer; every load-bearing step is inside the delivered span.\par\medskip
\begin{lemma}
\label{lemma-dualizing-schemes}
Let $K$ be a dualizing complex on a locally Noetherian scheme $X$.
Then $K$ is an object of $D_{\textit{Coh}}(\mathcal{O}_X)$
and $D = R\SheafHom_{\mathcal{O}_X}(-, K)$ induces an anti-equivalence
$$
D :
D_{\textit{Coh}}(\mathcal{O}_X)
\longrightarrow
D_{\textit{Coh}}(\mathcal{O}_X)
$$
which comes equipped with a canonical isomorphism
$\text{id} \to D \circ D$. If $X$ is quasi-compact, then
$D$ exchanges $D^+_{\textit{Coh}}(\mathcal{O}_X)$ and
$D^-_{\textit{Coh}}(\mathcal{O}_X)$ and induces an anti-equivalence
$D^b_{\textit{Coh}}(\mathcal{O}_X) \to D^b_{\textit{Coh}}(\mathcal{O}_X)$.
\end{lemma}

\begin{proof}
Let $U \subset X$ be an affine open. Say $U = \Spec(A)$ and
let $\omega_A^\bullet$ be a dualizing complex for $A$
corresponding to $K|_U$
as in Lemma \ref{lemma-equivalent-definitions}.
By Lemma \ref{lemma-affine-duality} the diagram
$$
\xymatrix{
D_{\textit{Coh}}(A) \ar[r] \ar[d]_{R\Hom_A(-, \omega_A^\bullet)} &
D_{\textit{Coh}}(\mathcal{O}_U) \ar[d]^{R\SheafHom_{\mathcal{O}_X}(-, K|_U)} \\
D_{\textit{Coh}}(A) \ar[r] &
D(\mathcal{O}_U)
}
$$
commutes. We conclude that $D$ sends $D_{\textit{Coh}}(\mathcal{O}_X)$ into
$D_{\textit{Coh}}(\mathcal{O}_X)$. Moreover, the canonical map
$$
L
\longrightarrow
R\SheafHom_{\mathcal{O}_X}(K, K) \otimes_{\mathcal{O}_X}^\mathbf{L} L
\longrightarrow
R\SheafHom_{\mathcal{O}_X}(R\SheafHom_{\mathcal{O}_X}(L, K), K)
$$
(using Cohomology, Lemma \ref{cohomology-lemma-internal-hom-evaluate}
for the second arrow)
is an isomorphism for all $L$ because this is true on affines by
Dualizing Complexes, Lemma \ref{dualizing-lemma-dualizing}\footnote{An
alternative is to first show that
$R\SheafHom_{\mathcal{O}_X}(K, K) = \mathcal{O}_X$ by
working affine locally and then use
Lemma \ref{lemma-internal-hom-evaluate-isom} part (2)
to see the map is an isomorphism.}
and we have already seen on affines that we recover what
happens in algebra.
The statement on boundedness properties of the functor $D$
in the quasi-compact case also follows from the corresponding
statements of Dualizing Complexes, Lemma \ref{dualizing-lemma-dualizing}.
\end{proof}

\end{document}
